\section{Evaluation}

Although a persistent name represents the history of a dart, an orbit is subject
to topological changes and, thus, requires the cons\-truc\-tion of its own history
in order to be accurately matched at reevaluation. %
Once the initial evaluation (Fig.~\ref{fig:evaluation}) is done and its
parametric specification (Fig.~\ref{fig:init-spec}) has been built, an evaluation's Directed
Acyclic Graph (or DAG) must be issued for each persistent name before
any parametric specification can be reevaluated. %
An e\-va\-lua\-tion DAG traces the history of each topological parameter back to the
first created orbits it originates from, thus allowing the matching of the
corresponding topological parameter at reevaluation time. %

\subsection{Evaluation DAG}

\begin{figure}[ht]
  \centering
  \includegraphics[width=0.25\textwidth]{pn3}
  \caption{$PN_{3}$'s evaluation DAG}
  \label{fig:eval-dag-pn3}
\end{figure}

An evaluation DAG is built parsing %with 
the applications and nodes
of a persistent name from end to start. %
It is sorted by levels representing the different applications
inside a history. %
Each level is made of an orbit level and an event level.
The orbit level contains a node's name and some orbits.
The event level contains a rule's application number and some events. %

For example, let us consider $PN_{3}=[1n_5;2n_6;3n_0]$ (Fig.~\ref{fig:4-eval}). 
% Début XS
$PN_3$ represents the topological parameter upon which the face extrusion rule is applied. This rule is filtered by hook $n_1$, matching the face $\orbit{0, 1}(47)$ (Fig.~\ref{fig:3-eval}).
% End XS
%c'est le paramètre topo de la règle d'extrusion appliquée sur la face 0~,~1~(~47~), filtrée par le hook n1.
%It is the topological parameter of the extrusion rule applied to face 0~,~1~(~47~), filtered by hook n1.
The matched orbit's history is built from its dart's history (\textit{i.e.} its persistent name). %
%\maxime{!} XS ??
The DAG is built bottom-up by a backward traversal through the persistent name. %
% Début XS
Since $PN_3$ is made of 3 parts, its evaluation DAG contains 3 levels.
The last part of $PN_3$ is $3n_0$, meaning that the dart of interest is filtered by the node $n_0$ of the 
third operation in the initial parametric specification.
Therefore, the Orbit level 3 contains both $n_0$ and the matched orbit $\orbit{0, 1}$.
As shown in (Fig.~\ref{fig:init-spec}), the third operation is the vertex in\-ser\-tion on an edge (Fig.~\ref{fig:rule-insert-vertex}). The right side of this rule has modified the orbit $\orbit{0,1}(n_0)$. We infer that the Event level 3 contains the third operation (\verb+3-insert+) and the event's name \texttt{MODIFICATION}.
% End XS
%Comme l'indique le pn3, le brin d'intérêt est filtré par le noeud n0 de la règle d'insertion (application 3) (Fig.~\ref{fig:rule-insert-vertex}). 
%As indicated by pn3, the dart of interest is filtered by node n0 of the insertion rule (application 3)
%le membre droit de la règle d'insertion a modifié l'orbite 01(n0).
%The right-hand side of the insertion rule has modified orbit 01(n0). 
%The third orbit level contains the node name $n_{0}$ and a volume face orbit $\orbit{0,1}$.
%The third event level contains the application name \texttt{3-insert} and the event's name \texttt{MODIFICATION}. %

%A syntactic analysis of the extrusion rule (Fig.~\ref{fig:rule-face-extrusion})for the face $\orbit{0,1}(n_{0})$ reveals that it was \emph{modified} by the vertex insertion rule's application $3$ so the contained event is a modification. %
Continuing through the persistent name's backward traversal, the previous element \texttt{2n6} allows determining the DAG's second level %computing 
in a similar way.
%with the extrusion rule's application $2$.
%The orbit level is indexed with $n_{6}$ and
%since the face orbit was modified (\textit{i.e.} no creation nor split nor
%merging of the orbit), it can be traced up to this level.
%The event level, this
%time, is indexed by $2$ and a query will find that the face orbit was created by
%this application. %
Finally, \texttt{1n5} allows computing the first level.   
%for the square creation rule's application $1$ (Fig.~\ref{fig:rule-create-square}):
%The orbit level is indexed with $n_{5}$ and
%since in the application $2$ the face $\orbit{0,1}$-orbit incident to $n_{6}$
%originates from another face, the current orbit level contains a
%$\orbit{0,1}$-orbit; The event level above is indexed by $1$ and a syntactic
%analysis of the square rule reveals that the face $\orbit{0,1}$-orbit was
%created. %
%The construction of the DAG finishes after application 1 since the left-hand
%side of the rule does not contain any node. %

At last, the produced evaluation DAG (Fig~\ref{fig:eval-dag-pn3}) represents the volume face $\orbit{0,1,3}$-orbit's history  resulting from applying \texttt{4-extrusion}.
This DAG can be read top-down:
\begin{description}
    \item[Level 1] The application \texttt{1-square} creates the volume face $\orbit{0,1}(n_5)$.
    \item[Level 2] From this volume face, \texttt{2-extrusion} creates the vo\-lu\-me face $\orbit{0,1}(n_6)$.
    \item[Level 3] Finally, this latter volume face is modified by \texttt{3-insert}, inserting a vertex on its edge.
\end{description}
%This DAG can be interpreted in a top-down fashion

%For example with the $PN_{3}$ which traces dart $47$'s history (Fig~\ref{fig:3-eval}: 

The %extrusion rule's application 
application of the extrusion rule
matches face $\orbit{0,1}(n_{5})$. Then, the ex\-tru\-sion creates the face $\orbit{0,1}(n_{6})$ (Fig.~\ref{fig:2-eval}).
Finally, %the vertex insertion rule's application
applying the vertex insertion rule
modifies the face $\orbit{0,1}(n_{0})$ (Fig.~\ref{fig:3-eval}). %

% Début XS
The syntactic analysis of the rules enables events to be computed only once. These events can be stored in cache to automatically build other evaluation DAGs.
% Fin XS
%The syntactic analysis of rules drives to the computation of events once. 
%Then it may be stored for later use in order to automatically build an evaluation DAG. %
%Then it is used many times to build evaluation DAGs. %

\subsection{Traces and origins}

In order to accurately represent the history of an orbit, two types of arrows are used in an evaluation DAG (and later in the reevaluation DAG): black trace arrows and red origin arrows. 

A \emph{black trace arrow} allows orbit evolution tracing. For example, the triangulation rule splits an initial face into multiple subfaces. If one of these subfaces is referenced in the DAG, it is connected by a trace arrow to the initial face. Therefore, a trace arrow connects two orbits of the same dimension.
A \emph{red origin arrow} allows lin\-king an orbit with the orbit that generated it, thereby connecting two orbits of potentially different dimensions. For example, continuing with the triangulation rule (Fig.~\ref{fig:rule-triangulation}), we observe that upon applying this rule, each edge of the initial face generates a different subface. If any of these subfaces is referenced in the DAG, it is then connected by an origin arrow to the edge that generated it.
This is what can be observed considering again the example in Fig.~\ref{fig:evaluation} and more pre\-ci\-se\-ly the colour rule's $PN_7$ parameter. $PN_7$ represents dart $35$'s history and $PN_7$'s evaluation DAG (Fig.~\ref{fig:eval-dag-pn7}) represents the history of the volume face that needs to be colored (the volume face ad\-ja\-cent to dart $35$).
This volume face $\orbit{0,1}(n_0)$ is the result of %the 5-triangulation rule 
operation \verb+5-triangulate+ 
which splits the volume face $\orbit{0,1}(n_4)$ and has the face edge $\orbit{0}(n_4)$ as its origin (respectively represented by a black and red arrow between orbit level 2 and event level 3).

As explained in section 4.1, this DAG is built using a  traversal of $PN_7$ and a bottom-up construction. This process is done in a similar way for both traces and origins, allowing for an efficient persistent naming mechanism that also takes into account the im\-pact of origin modifications during reevaluation.
To illustrate this, let us consider the previous face edge $\orbit{0}(n_4)$, which is the origin of the volume face $\orbit{0,1}(n_0)$ that needs to be colored. Suppose that, due to the addition of an operation in the edited specification, this origin %is eventually
may be 
split into two face edges. Upon applying %the 5-triangulation, 
\verb+5-triangulate+,
these two face edges will generate two volume faces, which can then be matched during reevaluation to the face to be colored, leveraging on the  origin orbit recorded in the DAG.

In a formal way, an origin orbit can be automatically deduced through the syntactic analysis of a rule.
More precisely, if $n$ is a hook and $n'$ is not a preserved node different from $n$, the origin of an orbit $R\orbit{o}(n')$ is the suborbit $L\orbit{o'}(n)$ consisting of the set of $n$'s implicit arcs which are:
\begin{itemize}
    \item rewritten on $R\orbit{o}(n')$;
    \item not rewritten on $R\orbit{o}(n')$ and belonging to $\orbit{o}$.
\end{itemize}

For example, let us calculate the origin of a volume face ($\orbit{0,1}$-orbit) resulting from a split in the triangulation rule (Fig.~\ref{fig:rule-triangulation}). Only the implicit arc $0$ in the $L\orbit{0,1}$-orbit incident to the hook $n_0$ is rewritten on $R\orbit{0,1}(n_0)$). Hence, the origin orbit of a volume face split by the triangulation rule is a face edge ($\orbit{0}$-orbit).

When the orbit is either \emph{split} or \emph{merged}, the
syntactic analysis al\-lows the deduction of an \emph{origin} in addition to a traced orbit.
When the orbit is \emph{created}, if the left-hand side of the rule is empty (mea\-ning the orbit is created from scratch), there is neither trace nor origin (event level 1 in Fig.~\ref{fig:eval-dag-pn7} where the 1-square rule creates the edge face $\orbit{0}(n_3)$ from scratch). Otherwise, there is no trace but an origin (event level 2 in Fig.~\ref{fig:eval-dag-pn7} where the 2-extrusion rule creates the front face $\orbit{0,1}(n_4)$ of the cube from the origin previous face edge $\orbit{0}(n_3)$).
When the orbit is just \emph{modified} or \emph{not modified}, there is only a traced orbit (event level 3 in Fig.~\ref{fig:eval-dag-pn3} where the 3-insert rule modifies the top face $\orbit{0,1}(n_6)$ of the cube inserting a vertex on its boundary).

\begin{figure}[hbt]
  \centering
  \includegraphics[width=0.3\textwidth]{pn7}
  \caption{$PN_{7}$'s evaluation DAG}
  \label{fig:eval-dag-pn7}
\end{figure}



\subsection{Paths to origins}

\begin{figure}[hbt]
  \centering
  \includegraphics[height=35pt]{rule-face-collapse}
  \caption{Face collapse rule}
  \label{fig:rule-face-collapse}
\end{figure}

%Nous venons de définir l'origine d'une orbite $R\orbit{o}(n')$ lorsque $n'$ n'est pas un noeud préservé différent du hook (cf def précédente). Dans le cas où il s'agit d'un noeud préservé différent du hook, il convient d'ajouter dans l'origine, le chemin permettant de rejoindre les liaisons implicites du hook depuis $n'$. En effet, les liaisons implicites d'une origine sont celles du hook et non celles de $n'$.
We have defined the origin of an orbit $R\orbit{o}(n')$ when $n'$ is not a preserved node different from the hook.
In the opposite case, it is necessary to add to the origin the \emph{path} that allows reaching the implicit arcs of the hook from $n'$, because the implicit arcs of an origin are those of the hook, not those of $n'$.

For example, in the face collapse rule (Fig.~\ref{fig:rule-face-collapse}), node $n_2$ is a preserved node different from the hook. Assume we want to define the origin of the volume vertex $\orbit{1,2}(n_2)$.
A path represents the traversal in the right-hand side of the rule from the node $n_{2}$ to the hook node $n_{0}$. The traversed explicit arcs from $n_{2}$ to $n_{0}$ are, in the following order, $1$ and $2$ (written $@1.@2$ in the evaluation DAG).
This can be seen in the $PN_{6}$'s evaluation DAG (Fig.~\ref{fig:eval-dag-pn6}), where  %the application of the 6-collapse rule 
applying \verb+6-collapse+ 
generates a merge of vertices and the origin of the volume vertex $\orbit{1,2}(n_{2})$ is the volume face $@1.@2.\orbit{0,1}(n_{4})$.
%Concretely, 
Actually,
node $n_4$ of the extrusion rule (Fig.~\ref{fig:rule-face-extrusion}) used at the pre\-vious level of the DAG (level 4) matches dart $55$ on the lateral faces of the cube (Fig.~\ref{fig:before-after-extrusion}) because it is the dart matching the his\-to\-ry stored in $PN_6$ DAG (\verb+1-square+ creates the initial bottom vo\-lu\-me face $\orbit{0,1}(n_5)$. Then, \verb+2-extrusion+ applied on this bottom face creates the volume face $\orbit{0,1}(n_6)$ which is then modified by \verb+3-insert+ in $\orbit{0,1}(n_0)$. Finally, \verb+4-extrusion+ applied on this modified volume face creates node $n_4$, which matched dart $55$).
The volume face $@1.@2.\orbit{0,1}(n_{4})$ reached starting from dart 55 and following links 1 and 2 is indeed the top face expected to be collapsed.    
\begin{figure}[hbt]
  \centering
  \includegraphics[width=0.3\textwidth]{pn6}
  \caption{$PN_{6}$'s evaluation DAG}
  \label{fig:eval-dag-pn6}
\end{figure}

%originates from the volume face $\orbit{0,1}(n0)$ incident to the hook. (illustrated if Fig.~\ref{fig:5-eval} and Fig.~\ref{fig:6-eval}).

%In some cases, usually in a merging event, the origin orbit can be deleted after the application of a rule and the origin does not represent the right orbit. %

%Now, considering the $PN_{6}$, whose application $6$ is the collapse rule, the
%origin to the vertex $\orbit{1,2,3}$-orbit, a face $\orbit{0,1,3}$-orbit, is
%expected to be incident to the application $4$'s node $n{4}$. %
%However, the face incident to $n{4}$ is adjacent to the face matched by $n{0}$
%(Fig.~\ref{fig:demo-fig-extrusion-right-match}) although the application is
%supposed to collapse the top face matched by $n{6}$ (Fig.~\ref{fig:6-eval}). %

%In order to circumvent such a limitation, a \emph{path to the origin} must be maintained in the history. % For example, in $PN_{6}$ evaluation DAG (Fig.~\ref{fig:eval-dag-pn6}), a path $@1.@2.$ is associated to the face $\orbit{0,1}$-orbit incident to the application $4$'s node $n{4}$. % Since an origin is always in the hook node, a path to the origin represent a traversal from the current node, $n{2}$ in this case, to the hook node. % As a result, the traversed explicit arcs from $n{2}$ to $n{0}$ are, in the following order, $1$ and $2$. % In the Jerboa notation, this path between is written $@1.@2$.   Concretely, this means that the origin orbit to the vertex resulting from the collapse is, for any dart matched application $4$'s node $n{4}$, is the face orbit incident to $n{4}.@1.@2.\orbit{0,1}$ being the $\orbit{0,1}$-incident to $n{6}$ being, this time, the top face expected to be collapsed (Fig.~\ref{fig:5-eval}). %


%%% Local Variables:
%%% mode: latex
%%% TeX-master: "../main"
%%% End:
